\section{Security arguments}

\subsection{Security argument for Theorem~\ref{thm:bb}}\label{sec:bb-sec}

We prove security of Theorem~\ref{thm:bb} via simulation.

Let $\adv$ be a polynomial-time adversary interacting with the our protocol,
denoted hereafter $\Pi$.  In the following, we construct a simulator $\simu$
which, based on its interactions with ideal functionality $\F{txid}$, simulates an
execution of $\Pi$ to $\adv$, and then prove that $\adv$ cannot distinguish
between a real execution of $\Pi$ and a simulated one.

More specifically, we define a sequence of games, starting from a real-world
execution of $\Pi$ and ending with ideal functionality $\F{txid}$, and demonstrate
that each game is (computationally) indistinguishable from the previous one.

In the following, $\simu$ maintains the following tables, which are analogous
to $\reg$, $\perm$, and $\txns$ maintained by $\F{txid}$:

\begin{enumerate}

  \item $\regsimu[\user] = \langle \sk, \utk, \signature\rangle$, logging
    registered users;

  \item $\permsimu[(\auditor,\user,\epoch)] = \ek$, logging tracing permissions
    and epoch keys;

  \item $\txnssimu[(\user,\epoch)]$, logging submitted transactions.

\end{enumerate}

\paragraph{Game 0.} This is the real execution of protocol $\Pi$.

\paragraph{Game 1.} This is an execution of $\Pi$ except that $\simu$ simulates
honest participants.  Namely, $\simu$ receives the inputs and generates the
outputs of the ledger, registration authority, honest users, and honest
auditors by following the instructions of $\Pi$.  By construction, Game 1 is
indistinguishable from Game 0.

\paragraph{Game 2.} This game is identical to Game~1, except that $\simu$ now,
instead of executing $\Pi$ during user registration, simulates the output of
the registration authority solely based on the output $\F{txid}$ upon $\IDregister$
calls.  For each $\IDregister$ request by some honest user, $\simu$ selects
random $(\sk, \utk)$ and simulates an execution of $\Pi$, and stores
$\regsimu[\user] = \langle\sk, \utk, \signature\rangle$.

Upon receiving a valid registration request from a corrupt user,
$\simu$ issues a $\IDregister$ request to $\F{txid}$ on behalf of the user and returns
the corresponding output.  If $\F{txid}$ successfully registers the user, then
$\simu$ uses the extractability property of proof to extract pair $(\sk, \utk)$
and records $\regsimu[\user] = \langle \sk, \utk, \signature\rangle$, where
$\signature$ is computed following $\Pi$'s instructions and sent to $\user$.

This game is indistinguishable from Game 1 unless $\F{txid}$'s output upon one of
the $\IDregister$ calls diverges from that of an execution of the registration
functionality of $\Pi$.  This only occurs if the registration authority, in the
real world, accepts a registration request that $\F{txid}$ rejects, or vice versa. 
 Due to the correctness of $\Pi$, the event that
$\F{txid}$ accepts while $\Pi$ rejects will never occurs.  Similarly, the event that
$\F{txid}$ rejects while $\Pi$ accepts will never occur; in fact, this event can
only take place if a corrupt user is able to trick the registration authority,
in the real world, to register them twice.  However, since the registration
authority keeps track of registration requests, this will never succeed.

\paragraph{Game 3.} This is identical to Game~3, except that $\simu$ simulates
the responses of honest users to tracing authorisation requests, based only on
the output of $\F{txid}$ upon $\IDgrant$ calls.

Upon a $\IDgrant$ call to functionality $\F{txid}$, for epoch $\epoch$, honest user
$\user$, and honest auditor $\auditor$, $\simu$ updates
$\permsimu[(\auditor,\user,\epoch)] \sample \epsilon$ and simulates the
transcript between $\user$ and $\auditor$ using randomly generated information.
Given that information between $\user$ and $\auditor$ is exchanged using a
secure channel, this is indistinguishable from an execution of $\Pi$, in the
real world.

Upon an authorisation request, in the real world, from a corrupt auditor
$\auditor$ intended for honest user $\user$, $\simu$ forwards the request to
$\user$ in the ideal world.  $\user$, accordingly, calls $\IDgrant$, if $\F{txid}$
accepts, then $\simu$ retrieves triple $\langle \sk, \utk, \signature\rangle
\sample \regsimu[\user]$, simulates $\user$'s response following $\Pi$'s
instructions, and sets $\permsimu[(\auditor,\user,\epoch)] \sample \FE(\utk,
\epoch)$.

It is easy to infer that the simulated response is indistinguishable from an
actual response generated by $\user$ during an execution of $\Pi$.  Therefore,
Game 3 is indistinguishable from Game 2.

\paragraph{Game 4.} This is identical to Game 4, except that upon an
authorisation request from a corrupt auditor for some honest user $\user$ and
epoch $\epoch$, $\simu$ randomly samples epoch key $\ek$, instead of computing
$\FE(\utk, \epoch)$, sets $\permsimu[(\auditor,\user,\epoch)] \sample \ek$, and
then simulates the proof for relation $\rel_1$.

Furthermore, when simulating an honest user submitting a transaction to
the ledger, $\simu$ produces the transaction following the description of
$\Pi$, except that since now $\ek$ is sampled randomly, it simulates the ZK
proof for relation $\rel_3$.  Now, thanks to the pseudo randomness of $\FE$, the
zero knowledge and simulatability properties of proofs, we conclude
that this game is indistinguishable from Game 3.

\paragraph{Game 5.} This is identical to Game~4, except that $\simu$ now also
simulates the output of transaction submission based on the output of $\F{txid}$
in response to $\IDsubmit$ requests from honest users.

For each call by an honest user to $\IDsubmit$ interface during epoch $\epoch$,
$\simu$ checks if there is a corrupt auditor with a tracing permission that can
trace the transaction back to its origin.  This is achieved by issuing $\IDidentify$
queries on behalf of corrupt auditors for all the honest users that granted
them tracing permissions, until there is a match.  More specifically, $\simu$
executes the following steps: 

\begin{enumerate}

 \item For each pair of honest
  user and corrupt auditor $\langle \user, \auditor\rangle$ with
  $\permsimu[(\auditor,\user,\epoch)]\neq \perp$, call $\IDidentify$ and receive
  $\txns \sample [(\user,\epoch)]$.  
  
  \item If $\txns[(\user,\epoch)] \neq
  \txnssimu[(\user, \epoch)]$, let $\algo{TX} = \txns[(\user,\epoch)] \setminus
  \txnssimu[(\user, \epoch)]$, update  $\txnssimu[(\user, \epoch)] \sample
  \txnssimu[(\user, \epoch)]$, and output pair $\langle \user,
  \algo{TX}\rangle$. 
  
   \end{enumerate} Note that since steps 1 and 2 are executed
  every time an honest user calls $\IDsubmit$ with transaction $\tx$, we have, by
  construction, for any honest user $\user$ who granted a tracing permission
  for epoch $\epoch$ to a corrupt auditor $\txns[(\user,\epoch)] =
  \txnssimu[(\user, \epoch)]$ or $\txns[(\user,\epoch)] =  \txnssimu[(\user,
  \epoch)] \cup \tx$.  Furthermore, we have either $\txns[(\user,\epoch)] =
  \txnssimu[(\user, \epoch)]$ for all such users,  or there exists one user
  $\user^\star$ such that $\txns[(\user^\star,\epoch)] =
  \txnssimu[(\user^{\star}, \epoch)] \cup \tx$ and that's the user who called
  $\IDsubmit$.  If $\simu$ encounters the first scenario, then this implies that
  the honest user who called $\IDsubmit$ did not grant tracing permissions to any
  of the corrupt auditors.  If it encounters the second, then it concludes that
  no corrupt auditor is authorised to trace that honest user's transactions.
  Accordingly, we consider the following two cases.

  \emph{Case 1. No corrupt auditor has a tracing permission for the honest user
  for the epoch.} $\simu$ sets the confidential payload to some arbitrary
  ciphertext and instead of computing the commitment as $\com = \ComCommit(\ck,
  \utk, \openinfo)$ and the tag as $\Tag_\gamma = \FT(\ek,\gamma)$, it samples
  uniformly at random commitment $\com^{\star}$ and tag $\Tag_{\gamma}^\star$
  and simulates the corresponding proof and ZK proof.  Given that epoch keys are
  generated at random, we can safely argue that epoch keys provided for
  preceding and succeeding epochs are independent from the current epoch.  As a
  result, we can use the hiding property of the commitment scheme, the pseudo randomness
  property of $\FT$, and the simulatability and zero knowledge properties of
  proofs to conclude that this is indistinguishable
  from submitting transactions during an execution of $\Pi$.

  \emph{Case 2. There is a corrupt auditor with valid tracing permission for
  the honest user.} $\simu$ in this case is able to get pair $\langle \user,
  \algo{TX} \rangle$ where $\user$ identifies the honest user and $\algo{TX}$
  consists of the transaction $\tx$ that $\user$ sent with their $\IDsubmit$
  query.  Given this information, $\simu$ retrieves tuple $\langle\sk,
  \utk,\signature \rangle \sample \regsimu[\user]$ and the random epoch key
  $\ek \sample \permsimu[(\auditor,\user,\epoch)]$.  $\simu$ then, in the real
  world, sets the confidential payload $\payload$ to $\tx$, and simulates the
  rest of the transaction similarly to Game 4.  \medbreak

  \paragraph{Game 6.} This is identical to Game~5, except that $\simu$ now
  simulates the output of transaction submission in response to transactions
  from corrupt users, based on the output of $\F{txid}$'s $\IDsubmit$ interface.

  When a corrupt user $\user$ submits a transaction with a valid tag (i.e., a
  tag with valid ZK proof and proof and which does not appear in the ledger),
  $\simu$ uses the extractability property of proofs and the knowledge soundness
  of ZK proofs to extract $\user$'s witness.  Furthermore, thanks to the
  existential unforgeability of the digital signature used during registration,
  $\simu$ will be able to identify $\user$ using table $\regsimu$.  Next,
  $\simu$ sends a $\IDsubmit$ request on behalf of $\user$ to $\F{txid}$ and returns
  the corresponding output.  If $\F{txid}$ accepts the transaction, then $\simu$
  records it in $\txns[(\user,\epoch)]$.

  This game is indistinguishable from Game 2 unless $\F{txid}$ rejects a
  transaction that $\Pi$ deems valid.  We recall that $\F{txid}$ rejects a
  transaction in one of two cases:

  \begin{enumerate}

    \item \emph{$\user$ is not registered:}  this cannot be the case given that
      each entry in $\regsimu$ corresponds to a successful $\IDregister$ call to
      $\F{txid}$.

    \item \emph{$\user$ already submitted $\Max$ transactions in the current
      epoch:} Given the knowledge soundness of ZK proofs and determinism of
      $\FT$, $\Pi$ will not accept a transaction from $\user$ in this case as
      the tag will already appear in the ledger.

  \end{enumerate}

  \paragraph{Game 7.} This is identical to Game~6, except that $\simu$ aborts
  the game if the output of the $\IDidentify$ interface differs from the output of the
  tracing functionality of protocol $\Pi$.

  Note that $\simu$ aborts if one of these two events occurs.

  \begin{itemize}

    \item abort 1. A coalition of corrupt auditors and users are able to trace
      a transaction submitted by an honest user during an epoch, even though
      the user did not grant permissions to any corrupt auditor for that epoch.

    \item  abort 2.  An honest auditor fails to trace a transaction  of a
      corrupt user that was submitted during an epoch, although they have been
      granted the authorization.

  \end{itemize}

  We now show that the probability of these events is negligible.

  abort 1. Since the transactions of honest users during an epoch for which the
  honest users did not grant tracing permissions to any corrupt auditor are
  generated at random, this event does not occur.

  abort 2. We recall that a transactions $\tx$ from a corrupt user is only
  recorded in  $\txnssimu$ and forwarded to $\F{txid}$, if it carries a valid tag
  $\Tag_{\gamma}$; that is, the proof for relation $\rel_2$ and the ZK proof for
  relation $\rel_3$ are both valid.  Given the extractability property of proofs
  and knowledge soundness of ZK proofs, we can safely conclude that
  $\Tag_{\gamma} = \FT(\ek, \gamma) \ \wedge \ \ek = \FE(\utk, \epoch)$ with
  $\utk$ verifying Equations \ref{eq:rel2} and \ref{eq:rel3}.  Therefore, an
  honest auditor with the right permission will fail to trace $\tx$ only if
  they have accepted a bogus epoch key $\ek^\star$.  However, given the
  extractability property of the proof used to prove $\rel_1$ and the existential
  unforgeability of digital signatures, this event occurs only with a
  negligible probability.

  \subsection{PRF security arguments}

  \subsubsection*{PRF Security of $\FE$}\label{sec:fe-prf-proof}

  \begin{lemma}\label{lem:fe-prf} Under the DDH assumption in $\group{1}$ and
  modeling $\HE$ as a random oracle, the function $\FE(\utk, \epoch) =
\HE(\epoch)^{\utk}$ is a secure pseudorandom function.  \end{lemma}

We show that no polynomial-time adversary can distinguish $\FE(\cdot, \cdot)$
from a random function with non-negligible advantage.

Let $\adv$ be an adversary with oracle access to either $\FE(\utk, \cdot)$ or a
random function $R: \{0,1\}^* \to \group{1}$. Let $\adv$ distinguish $\FE(\utk,
\cdot)$ from a random function with advantage $\varepsilon > 0$. The advantage
is defined as $\varepsilon = | Pr[\adv = 1 | \FE] - Pr[\adv = 1 | R]|$.  We use
$\adv$ to construct a simulator $\simu$ that breaks DDH in $\group{1}$.

$\simu$ receives a DDH challenge $(\gen, \gen^a, \gen^b, Z)$ where $Z =
\gen^{ab}$ (DDH) or $Z \sample \group{1}$ (random). We use the distinguisher
that can tell apart values output by $\FE$  and random values to build an
adversary who can distinguish DDH tuples from random challenges.  It does the
following:

\begin{itemize}

  \item Implicitly, the secret key is $\utk^* = b$.

  \item For queries $\epoch_i$ to $\HE(\epoch_i)$, set $\HE(\epoch_{i}) =
    \gen^{a \cdot {r_i}} = (\gen^a)^{r_i} $ for fresh random $r_i \sample \Zp$.

  \item For queries $\epoch_i$ to $\FE(\utk^*, \cdot)$, set $\FE(\utk^*,
    \epoch_{i}) = Z^{r_i}$.

\end{itemize}

If $Z = \gen^{ab}$, then the simulator has constructed $\FE$ with key encoded
in $\gen^b$ and responses are consistent across repeated queries.

If $Z \sample \group{1}$, then each response is independent and uniformly
random in $\group{1}$, since $Z$ is independent of $\gen^a$ and $\gen^b$.

An adversary that distinguishes $\FE$ from random can then distinguish between
the case where $Z$ is a valid DDH tuple (and the responses are outputs of
$\FE$) and the case where $Z$ is random (and responses are uniformly random).

We then have the following advantage: 
%
    \begin{equation*}
%
\Adv[\adv_{\PRF}(\FE)] = \Adv[\adv_\mathsf{DDH}(\simu)].
%
    \end{equation*}
%

\subsubsection*{PRF Security of $\FT$}\label{sec:ft-prf-proof}

\begin{lemma}\label{lem:ft-prf} Under the BDDH assumption in bilinear groups,
the function defined as follows:
%
    \begin{equation*}
%
\FT(\ek, \gamma) = e(\ek, \prod_{i=0}^\ell
G_i^{x_i}),
%
    \end{equation*}
%
 is a secure pseudorandom function for well-formed inputs
$\gamma$.  \end{lemma}

Given BDDH challenge $(\gen, \gen^a, \gen^b, \ggen, \ggen^c, Z)$ where $Z =
e(\gen, \ggen)^{abc}$ (BDDH tuple) or $Z \sample \group{T}$ (random), an adversary
whose goal is to break BDDH, proceeds as follows. 

 \begin{itemize}

  \item Implicitly, set the epoch key to $\ek^* = \gen^a$.

  \item Choose random values $\alpha_0, \ldots, \alpha_{\ell}, \beta_0, \ldots,
    \beta_{\ell} \sample \Zp$ and set $G_i$ as follows for for $i = 0, \ldots,
    \ell$:
%
    \begin{equation*}
%
      G_i = \ggen^{\alpha_i c + \beta_i} = (\ggen^c)^{\alpha_i} \cdot
      \ggen^{\beta_i}.
%
    \end{equation*}
%

  \item For queries $\FT(\ek, \gamma)$ with bit decomposition $(\gamma_0,
    \ldots, \gamma_\ell)$:

    \begin{itemize}

      \item Compute $\signature = \sum_{i=0}^\ell \alpha_i \gamma_i$ and $\tau
        = \sum_{i=0}^\ell \beta_i \gamma_i$.

      \item Return $Z^{\signature} \cdot e(\gen^a, \ggen)^{\tau}$.

    \end{itemize}

\end{itemize}

If $Z = e(\gen, \ggen)^{abc}$, then for each query $\gamma$, the following
holds:
%
\begin{equation*}
%
  \prod_{i=0}^\ell G_i^{x_i} = \ggen^{c\signature_j + \tau_j} \ \ ; \ \
  e(\gen^a, \prod_{i=0}^\ell G_i^{x_i}) = e(\gen^a, \ggen)^{c\signature_j +
  \tau_j} = Z^{\signature_j} \cdot e(\gen^a, \ggen)^{\tau_j}.
%
\end{equation*}
%
Therefore, the responses of the adversary are correctly distributed as the
output of $\FT$. If the distinguisher against $\FT$ signals that $\FT$ is a
PRF, we can conclude that the BDDH challenge was indeed a BDDH tuple.

If $Z \sample \group{T}$, each response $Z^{\signature} \cdot e(\gen^a,
\ggen)^{\tau}$ uniformly random in $\group{T}$, since $Z$, $\alpha_i$, $\beta_i$ are
uniformly random and independent of $\ek$ and other chosen values.  If the
distinguisher returns that $\FT$ is not a PRF, we can then conclude that the
BDDH challenge was not a BDDH tuple.

We then have the following advantage:
%
\begin{equation*}
%
  \Adv[\adv_{\PRF}(\FT)] = \Adv[\adv_\mathsf{BDDH}(\simu)].
%
\end{equation*}
%

\subsubsection*{$\rel_3$ and $\widehat{\rel}_3$ equivalence proofs}

Correctness and soundness are shown as follows.  With a witness satisfying
$\rel_3$, pick arbitrary $s, r$ and define $\Gamma, U, V$.  All six clauses of
$\widehat{\rel}_3$ follow directly.  With a witness satisfying
$\widehat{\rel}_3$, it is possible to start from the pairing equation $e(V,
\Gamma) = \Tag_\gamma^s \cdot e(V, H)^r$ to recover the original $\rel_3$
constraint $e(\ek, \prod G_i^{\gamma_i}) = \Tag_\gamma$ via injectivity of the
pairing.  \rebekah{@Rebekah: recall the equations so this section is
self-contained.  One should not assume that the reader will go back and forth
between the sections.}

Zero knowledge is shown as follows.  For $\rel_3$, we have the following public
information,
%
\begin{equation*}
%
  \inst_3 = (\HE, \{\gggen{G}_i\}_{i\in[1,\ell]}, \gen_2, Q, \com, \Max,
  \epoch, \Tag_\gamma).
%
\end{equation*}
%
We define $\widehat{\inst}_3$ as $(\inst_3, \gggen{H}, U, V)$.

$(U, V, \Gamma)$ are computed using ephemeral randomness $s, r \sample \Zp$ as
follows:
%
\begin{equation*}
%
  U = \HE(\epoch)^{s}, \quad V = \ek^{s}, \quad \Gamma =
  \gggen{H}^r\prod_{i=1}^{\ell} \gggen{G}_i^{\gamma_i}.
%
\end{equation*}
%
Given this, recovery of $(\utk, \{\gamma_i\}_{i \in [1,\ell]})$ would require
solving $V = U^{\utk}$ with $U = \HE(\epoch)^s$ for unknown $s$.

Let $\gen = \HE(\epoch)$, so $U = \gen^s$. The verifier observes $(\gen, U,
\ek, V)$ where $V = \ek^s$.  By the DDH assumption, the distribution of $(\gen,
\gen^s, \ek, \ek^s)$ is computationally indistinguishable from $(\gen, \gen^s,
\ek, \ek^t)$ for uniformly random $s, t \sample \Zp$. Therefore any adversary
with advantage $\epsilon$ in computing $\utk$ from $(U, V)$ would contradict
CDH, and distinguishing would contradict DDH.

$s, r$ are uniformly random and used only once per proof . The pairing relation
as follows:
%
\begin{equation*}
%
  e(V, \Gamma) = \Tag_\gamma^s \cdot e(V, \gggen{H})^r,
%
\end{equation*}
%
holds by construction, but knowledge of this does not improve the adversary's
advantage in recovering the witness.

Given a witness $\widehat{\witness}_3 = (\witness_3, r, s)$ satisfying
$\widehat{\rel}_3(\widehat{\inst}_3, \widehat{\witness}_3) = 1$, we recover the
original constraint of $\rel_3$.

Since the randomness $s, r \sample \Zp$ is freshly sampled for each proof
instance and independent of the witness, the simulator can generate $(U, V,
\Gamma)$ without knowledge of $(\utk, \openinfo, \{\gamma_i\})$. Therefore zero
knowledge of $\rel_3$ follows from honest-verifier zero knowledge of the Sigma
protocol for $\widehat{\rel}_3$.

\rebekah{The way you go about this is that given the statement for $\rel_3$, you
  should be able to construct a proof as depicted in the description.  $U, V$
  you can select them totally at random, $\Gamma = H^r\prod_{i=1}^l
  G_i^{\gamma_i}$. In theory, $\Gamma$ can also be selected at random but in
  that case we cannot use range proofs a black box.  Then one can use the
simulatability of Schnorr proofs to show that one can prove all the statements
if one knows the challenge beforehand. It may make sense to say that we show
this for the interactive version, random oracle programmability can be used to
show it for the non-interactive version.}
